\chapter{Electroconvulsive Therapy}
\index{Electroconvulsive Therapy}

\section*{Definition and Overview}
Electroconvulsive therapy, known as ECT, is a treatment for severe mental illness in which a carefully controlled electrical stimulus is passed through the brain under general anaesthesia to produce a brief, therapeutic seizure. It is one of the most effective treatments for severe depression, particularly when medicines have not worked or when rapid improvement is needed because of risk to life. ECT is also used for some cases of severe mania and catatonia. The treatment is given in a course over several weeks, and while it carries a risk of memory effects, modern ECT is safe and can be life-saving.

\section*{Epidemiology}
ECT is used in a small proportion of people with depression, estimated at a few thousand people each year. It is more common in older people, in whom severe depression is common and antidepressant response may be slower, and in hospital settings where depression is severe enough to threaten life or require urgent treatment. Its use has declined since the advent of modern antidepressants, but it remains an important treatment for the most severe illness.

\section*{Aetiology and Pathophysiology}
The exact mechanism of ECT is not fully understood, but research points to several effects on the brain.
\begin{itemize}
    \item \textbf{Neurotransmitters:} ECT alters the levels and activity of brain chemicals involved in mood
    \item \textbf{Neuroplasticity:} the treatment promotes the growth of new connections between brain cells and increases levels of growth factors
    \item \textbf{Brain activity:} ECT changes patterns of brain activity in regions involved in depression
    \item \textbf{The seizure:} the therapeutic effect comes from a generalised seizure under anaesthesia, lasting about 30 to 60 seconds
    \item \textbf{Why it works:} the mechanisms overlap with those of antidepressants but act more quickly and powerfully, which is why ECT can work when medicines have failed
\end{itemize}

\section*{Clinical Features}
ECT is considered in specific clinical situations.
\begin{itemize}
    \item Severe depression with psychotic features, such as delusions or hallucinations
    \item Depression with a high risk of suicide, where rapid response is critical
    \item Depression that has not responded to multiple medication trials
    \item Severe mania, where ECT can act faster than medicines
    \item Catatonia, a state of severe motor disturbance, which responds dramatically to ECT
    \item Depression during pregnancy, where ECT may be safer than some medicines
\end{itemize}

\section*{Differential Diagnosis}
ECT is one option among several for severe illness.
\begin{itemize}
    \item Antidepressant medication, alone or in combination
    \item Transcranial magnetic stimulation, a gentler but less powerful brain stimulation treatment
    \item Psychological therapies, which are not usually sufficient for severe depression alone
    \item The decision between treatments is made by the psychiatrist with the patient
\end{itemize}

\section*{When to Seek Medical Care}
People experiencing severe depression with thoughts of self-harm, psychosis, or inability to care for themselves need urgent psychiatric assessment, where ECT may be discussed.

\textbf{Contact your local emergency services for:}
\begin{itemize}
    \item Thoughts or plans of suicide
    \item Refusing food and fluids or severe neglect of self
    \item Agitation, confusion, or psychosis that is dangerous
    \item Sudden collapse or loss of consciousness
\end{itemize}

\section*{Diagnosis and Investigations}
Before ECT, the person is fully assessed.
\begin{itemize}
    \item \textbf{Psychiatric assessment:} confirmation of the diagnosis and the indication for ECT
    \item \textbf{Medical assessment:} general health review, including the heart
    \item \textbf{Blood tests:} routine pre-anaesthesia checks
    \item \textbf{ECG and sometimes imaging:} to ensure safe anaesthesia
    \item \textbf{Consent:} a detailed discussion of the procedure, benefits, and risks, with consent required before treatment
\end{itemize}

\section*{Management}
ECT is delivered as a course of treatments.
\begin{itemize}
    \item \textbf{Anaesthesia:} a short general anaesthetic and a muscle relaxant are given before each treatment
    \item \textbf{The treatment:} electrodes on the scalp deliver a brief stimulus; the seizure is monitored
    \item \textbf{Course:} typically six to twelve treatments, two or three times a week
    \item \textbf{Assessment:} mood and side effects are reviewed throughout
    \item \textbf{Continuation:} maintenance ECT, antidepressant medicines, or both are used to prevent relapse
    \item \textbf{Outpatient ECT:} possible once the person is stable
\end{itemize}

\section*{Complications}
\begin{itemize}
    \item Memory loss, the most common concern: confusion and memory gaps during the course, which usually improve but can persist for some people
    \item Headache, nausea, and muscle aches after treatments
    \item Effects of anaesthesia
    \item Rarely, cardiac or other medical complications
    \item Relapse after the course without continuation treatment
\end{itemize}

\section*{Prognosis and Outcomes}
ECT is one of the most effective treatments in psychiatry, with remission rates of 60 to 80 per cent for severe depression, and it often works when medicines have failed. Improvement can occur within days. The main drawback is the risk of memory effects, which is discussed openly and balanced against the risks of untreated severe depression, including death by suicide. With continuation treatment, many people remain well.

\section*{Prevention and Self-Care}
\begin{itemize}
    \item Ask questions and discuss concerns about ECT fully with the treating team
    \item Bring a support person to appointments and assessments
    \item Plan practical support for memory effects during the course
    \item Take continuation treatment as recommended after ECT
    \item Attend follow-up psychiatric reviews
    \item Look after sleep, nutrition, and gentle activity during recovery
    \item Seek immediate help if suicidal thoughts return
\end{itemize}

\section*{Sources and Further Reading}
The following reputable sources provide further information for healthcare students and patients seeking reliable, current advice about electroconvulsive therapy.
\begin{itemize}
    \item Royal Australian and New Zealand College of Psychiatrists. \textit{ECT information for patients.} https://www.ranzcp.org/
    \item Beyond Blue. \textit{Depression treatments: ECT.} https://www.beyondblue.org.au/
    \item Healthdirect Australia. \textit{Electroconvulsive therapy.} https://www.healthdirect.gov.au/electroconvulsive-therapy-ect
    \item NHS. \textit{Electroconvulsive therapy (ECT).} https://www.nhs.uk/mental-health/treatments-and-wellbeing/electroconvulsive-therapy-ect/
    \item Mayo Clinic. \textit{Electroconvulsive therapy (ECT).} https://www.mayoclinic.org/tests-procedures/electroconvulsive-therapy/
    \item Head to Health. \textit{Treatments for severe depression.} https://www.headtohealth.gov.au/
\end{itemize}
